\chapter{Il sistema proposto}
\label{ch:core}

% prendere ispirazione dalla parte di formalization di questo articolo: https://iris.unica.it/bitstream/11584/461825/1/1-s2.0-S0164121225003516-main_2.pdf
% possibilmente pensare a una figura con uno schema a blocchi
% descrivere il sistema in modo generale, senza scendere nei dettagli implementativi

% NOTA: l'articolo indicato e' \cite{tony2025prompt} (Journal of Systems and
% Software, 2025): usarne la sezione di formalizzazione come modello per la
% notazione di questo capitolo.
%
% Schema a blocchi suggerito:
%   dataset (funzione) -> costruzione del prompt -> modello -> normalizzazione
%   dell'output -> file di test -> esecuzione -> misure
%
% Il capitolo descrive COME si misura. I numeri stanno nel capitolo 4.


\section{Formalizzazione}
\label{sec:formalizzazione}
% Notazione da introdurre, sul modello di \cite{tony2025prompt}:
% - F = insieme delle funzioni sotto test, f in F, |F| = 100;
% - M = insieme dei modelli considerati, m in M, |M| = 3;
% - P(f) = prompt costruito a partire dalla funzione f tramite il template fisso;
% - T(m, f) = file di test prodotto dal modello m per la funzione f;
% - t in T(m, f) = singola funzione di test contenuta nel file.
%
% La distinzione fra T(m,f) e t e' sostanziale e va introdotta qui: le metriche
% di esito si definiscono sul file, quelle di correttezza sul singolo test.
% Serve nel capitolo 4, dove i file che passano interamente sono pochissimi
% mentre i test corretti sono abbastanza da permettere una media.
%
% Definire poi le metriche come funzioni su T(m,f) e f, distinguendo:
% - misure statiche: calcolabili senza eseguire nulla;
% - misure dinamiche: richiedono l'esecuzione dei test.


\section{Architettura generale}
% Il sistema e' una catena a colpo singolo: una funzione, una richiesta, un file
% di test. Non c'e' iterazione di riparazione, per scelta: si vuole misurare la
% capacita' dei modelli senza aiuti esterni. Va detto esplicitamente, perche' e'
% cio' che distingue questo impianto da CoverUp \cite{pizzorno2024coverup} e
% CodaMOSA \cite{lemieux2023codamosa} discussi nella Sezione~\ref{sec:llm}.
%
% Blocchi:
% 1. selezione del campione dal dataset;
% 2. costruzione del prompt secondo un template fisso;
% 3. invocazione del modello a parametri costanti;
% 4. normalizzazione della risposta;
% 5. salvataggio in un file separato per ogni coppia (modello, funzione);
% 6. esecuzione dei test in ambiente isolato e calcolo delle metriche.
%
% Sul blocco 4 va fatta una precisazione, perche' e' una scelta con un costo:
% la normalizzazione rimuove gli eventuali blocchi markdown dalla risposta prima
% di salvarla. Rende i file eseguibili, ma cancella la prova, e quindi impedisce
% di misurare a posteriori quante volte il vincolo "niente markdown" e' stato
% violato. Dichiararlo come limite noto e rimandare alla sezione sulle minacce
% alla validita'.


\section{Dataset e selezione dei campioni}
% ULT_Lite \cite{huang2025ult}: 200 campioni in formato JSON Lines; ne vengono
% usati i primi 100 nell'ordine del file (criterio deterministico e riproducibile).
% Ogni campione contiene nome della funzione, codice sorgente, descrizione in
% linguaggio naturale, identificativo e una lista di assert.
% Le funzioni vanno da 8 a 171 righe, con mediana 34.
%
% Due esclusioni da motivare:
% - la descrizione in linguaggio naturale non viene passata al modello, per
%   scelta sperimentale: si misura cosa il modello ricava dal solo codice;
% - la lista di assert non e' utilizzabile come riferimento, perche' gli autori
%   del benchmark dichiarano di non rilasciare i test di verita' per non
%   contaminare i futuri addestramenti. Verificato sui dati: quegli assert
%   passano solo nel 12% dei casi sul codice fornito.


\section{Costruzione del prompt}
% Template fisso in tre sezioni:
% [instruction] i vincoli: stile pytest senza classi, import esplicito della
%   funzione sotto test, da 3 a 8 funzioni di test con nomi numerati, almeno un
%   assert per test, divieto di riscrivere la funzione, divieto di spiegazioni e
%   markdown;
% [data] il solo codice della funzione;
% [format] lo scheletro atteso dell'output.
%
% Riportare il template per intero in un ambiente verbatim: e' parte del
% contributo e permette la riproducibilita'.
% Motivazione della sezione [format]: per i modelli piccoli vedere la forma
% attesa e' piu' efficace che leggerne la descrizione.
%
% Anticipare che i sei vincoli non sono decorativi: cinque di essi diventano la
% metrica di aderenza al formato del capitolo 4, e il loro rispetto e' cio' che
% separa nettamente il modello da 1B dagli altri due.


\section{Modelli e parametri}
% Tre modelli della stessa famiglia \cite{llama3herd} di dimensione crescente
% (1B, 3B, 8B), a parita' di prompt e parametri: temperatura 0 (generazione
% deterministica, una sola esecuzione per campione) e limite di 1024 token in
% uscita.
%
% ATTENZIONE: non scrivere che il limite di 1024 token basta sempre a chi
% rispetta il vincolo sul numero di test. I dati lo smentiscono: fra i file
% troncati ce ne sono con 3, 5, 6 e 7 funzioni di test, cioe' dentro il vincolo.
% Il budget dipende anche da quanto sono verbosi i singoli test, non solo da
% quanti sono. Il troncamento va quindi trattato come fenomeno a se'.
%
% Infrastruttura: esecuzione in locale con Ollama \cite{ollama} attraverso
% un'interfaccia compatibile con la libreria openai. I pesi sono quelli
% rilasciati da Meta, qui quantizzati a 4 bit: la quantizzazione riduce la
% precisione numerica dei parametri e semmai penalizza il modello, non lo
% favorisce. Va dichiarata.
%
% I primi 100 campioni dell'8B erano stati generati in precedenza su un servizio
% remoto: quei dati sono archiviati e usati solo come confronto secondario fra
% piena precisione e quantizzazione. Non fanno parte dello studio.


\section{Classificazione degli esiti}
\label{sec:albero}
% Sezione nuova, ed e' la struttura portante dei capitoli 3 e 4.
% Va accompagnata da una figura ad albero, qui vuota; il capitolo 4 la ripopola
% con le percentuali.
%
% N campioni
% |- test NON scritti .............. fallimento della pipeline (parte statica)
% \- test SCRITTI
%    |- sintatticamente non validi .. ast.parse fallisce (parte statica)
%    \- sintatticamente validi
%       |- non eseguibili ........... pytest non raccoglie test, import fallito
%       |- timeout .................. oltre il limite di tempo
%       |- falliti .................. almeno un test non passa
%       \- passati .................. tutti i test passano
%
% Il confine fra parte statica e parte dinamica cade sotto "sintatticamente
% validi": tutto cio' che sta sopra si stabilisce senza eseguire il codice.
%
% Ogni ramo si riporta in percentuale e separatamente per ogni modello.
% Un ramo puo' risultare vuoto: va comunque riportato, perche' la classificazione
% e' definita a priori e riportarla per intero e' cio' che la rende credibile.


\section{Esecuzione e misura}
\label{sec:esecuzione}
% Per ogni file di test generato:
% 1. controllo sintattico con il parser AST, senza eseguire il codice;
% 2. esecuzione con pytest \cite{pytest} in una cartella temporanea isolata e con
%    limite di tempo, accanto al file contenente la funzione sotto test;
% 3. calcolo della copertura con coverage.py \cite{coveragepy} sulla sola
%    funzione sotto test, di riga e di ramo.
%
% L'output integrale di pytest viene conservato per ogni campione: e' la base
% dell'analisi dei fallimenti descritta piu' avanti.
%
% --- misure statiche ---
% Verificate sull'albero sintattico, senza eseguire nulla: validita' sintattica,
% numero di funzioni di test, e i cinque vincoli del prompt controllabili
% (import corretto, numero di test nel range, nomi conformi, un assert per test,
% assenza di classi). Va aggiunta la variante che accetta come oracolo anche i
% blocchi pytest.raises: alla lettera violano il vincolo sull'assert, ma un
% oracolo ce l'hanno, e le due letture misurano cose diverse.
%
% --- correttezza ---
% Pass@1 definito sul SINGOLO test e non sul file, coerentemente con
% \cite{huang2025ult}: proporzione di funzioni di test che passano sul totale
% generato. La ragione e' anche pratica e va detta: i file in cui tutti i test
% passano sono pochissimi, e una metrica definita su quella base non avrebbe
% campioni sufficienti.
%
% --- copertura: eseguita contro verificata ---
% E' il punto metodologicamente piu' rilevante del capitolo.
% La copertura si calcola due volte:
% - con tutti i test del file: dice quanto codice viene ESEGUITO. Include
%   l'esecuzione prodotta dai test che falliscono, perche' un assert sbagliato
%   esegue comunque la funzione prima di sollevare l'errore;
% - con i soli test che passano, deselezionando i falliti e rieseguendo: dice
%   quanto codice viene davvero VERIFICATO da un'asserzione corretta.
% La differenza fra le due misura quanto la copertura sopravvaluta cio' che i
% test garantiscono. E' la forma operativa dell'argomento di Inozemtseva e
% Holmes \cite{inozemtseva2014coverage} discusso nella Sezione~\ref{sec:limiti}.
%
% Dichiarare sempre il denominatore delle medie: tutti i campioni, i soli
% eseguibili, oppure i soli campioni che hanno importato la funzione. La terza
% base serve a separare l'incapacita' di scrivere test dal non aver seguito
% l'istruzione sull'import.
%
% --- duplicazione ---
% Definizione operativa: quota di righe ripetute sul totale delle righe di test,
% ignorando righe vuote e commenti; una riga presente n volte contribuisce n-1
% ripetizioni. In percentuale, per essere indipendente dalla lunghezza della
% suite. Calcolata su due basi: tutte le righe di test, e le sole righe dei test
% che passano.
%
% --- analisi dei fallimenti ---
% Dai report di pytest si estrae, per ogni test fallito, il tipo di eccezione e
% il file in cui e' stata sollevata. Da qui le due famiglie:
% - errore di oracolo: il test esegue la funzione ma sbaglia l'aspettativa
%   (assert fallito, eccezione attesa e non arrivata, oppure eccezione sollevata
%   dentro la funzione perche' l'input non e' accettato);
% - errore d'uso: il test non arriva a esercitare la funzione.
% E' la traduzione empirica della distinzione fra raggiungere il codice e sapere
% cosa deve restituire, introdotta nella Sezione~\ref{sec:classici}.
%
% --- nota implementativa ---
% Il campo con il codice del dataset non include gli import necessari alla
% funzione, che vanno anteposti ed esclusi dal calcolo della copertura tramite
% "pragma: no cover". Senza questa correzione alcune funzioni sollevano NameError
% e i test risultano falliti per un motivo estraneo al modello.


\section{Minacce alla validità}
\label{sec:minacce}
% Sezione nuova. Tre punti emersi durante il lavoro, tutti da dichiarare.
%
% 1. Codifica dei caratteri. Tre funzioni del dataset contengono caratteri non
%    ASCII (giapponese in due casi, una freccia nel terzo). I modelli li
%    riproducono nelle asserzioni, e la stampa della risposta sollevava
%    UnicodeEncodeError sulla console di Windows, facendo fallire la
%    generazione. Corretto forzando UTF-8. Senza la correzione quelle tre
%    funzioni sarebbero state escluse dal confronto per una ragione che non
%    riguarda i modelli: e' un esempio di come un dettaglio infrastrutturale
%    possa falsare una misura.
%
% 2. Determinismo. Alcune funzioni del dataset usano numeri pseudocasuali o
%    dipendono dall'ordine di iterazione, quindi l'esecuzione dei test non e'
%    del tutto riproducibile: fra due esecuzioni successive le medie oscillano
%    di circa due decimi di punto. I valori vanno riportati con una cifra
%    decimale e non letti come esatti al decimo.
%
% 3. Quantizzazione. I modelli locali sono quantizzati a 4 bit. La quantizzazione
%    tende a peggiorare leggermente le prestazioni, quindi i risultati vanno
%    letti come una stima prudente e non ottimistica.
%
% 4. Generalizzabilita'. Cento funzioni di un solo benchmark, una sola famiglia
%    di modelli, una sola formulazione del prompt. Il ruolo del prompt e' noto
%    per essere rilevante \cite{tony2025prompt}: risultati diversi con un
%    template diverso non sarebbero sorprendenti.